\section{Seleksi Fitur (\textit{Feature Selection})}

Sesuai dengan pembahasan pada Subbab~\ref{sec:feature_selection}, tahapan pemilihan atribut terbaik (\textit{feature selection}) diimplementasikan untuk mengurangi beban komputasi algoritma dan mengabaikan informasi yang tidak relevan (\textit{noise}). Evaluasi fitur ini dijalankan setelah himpunan data dipisahkan menjadi data latih dan data uji (\textit{train-test split}). Pendekatan ini bertujuan untuk menutup celah kebocoran informasi dari data uji ke dalam proses evaluasi fitur (\textit{data leakage}). Berikut metode yang diterapkan.
\begin{enumerate}
    \item \textbf{Uji Chi-Square (\textit{SelectKBest}):} mengevaluasi seberapa kuat dependensi antara setiap variabel kategori dengan target klasifikasi secara independen.
    \item \textbf{Uji \textit{Mutual Information}:} mengukur seberapa besar suatu fitur mampu memberikan petunjuk tambahan atau mengurangi tingkat ketidakpastian sistem dalam memprediksi kelas penyewa.
    \item \textbf{Korelasi Pearson:} mengukur kekuatan hubungan garis lurus (linear) antara fitur numerik dengan target klasifikasi, yang mana tingkat signifikansinya dinilai berdasarkan nilai absolut (\textit{absolute value}) dari koefisien korelasi.
\end{enumerate}

Untuk mengintegrasikan hasil dari tiga metrik tersebut, setiap nilai uji dinormalisasi menjadi rentang standar antara 0 hingga 1 menggunakan fungsi algoritma \texttt{MinMaxScaler}. Rata-rata dari ketiga nilai yang telah dinormalisasi kemudian dikalkulasi untuk menghasilkan satu skor akhir terpadu (\textit{Total Score}). Berdasarkan hasil eksekusi program (\textit{output}) pada tahap seleksi fitur untuk data riil, sistem melakukan pemeringkatan seluruh variabel berdasarkan skor total tertinggi, yaitu \texttt{lama\_menginap}, \texttt{checkout\_hour}, dan \texttt{checkin\_hour}. Ketiga fitur terpilih inilah yang kemudian dipakai untuk ke tahap pelatihan algoritma agar model dapat bekerja dengan tingkat komputasi yang lebih efisien serta terhindar dari bias informasi yang tidak relevan.